\begin{table}[H]
\centering
\caption{Productividad laboral por departamento, 2014--2024pr}
\label{tab:pib_geih_productividad_departamento}
\scriptsize
\begin{tabular}{lrrrr}
\toprule
Departamento & PIB/trab. 2024 & PIB/hora 2024 & Crec. PIB/trab. & Crec. PIB/hora \\
\midrule
Quindío & 33,8 & 15,1 & 7,21\% & 8,21\% \\
Putumayo & 443,0 & 202,1 & 6,27\% & 6,62\% \\
Risaralda & 38,7 & 16,4 & 3,37\% & 4,14\% \\
Caquetá & 24,9 & 10,5 & 3,51\% & 3,88\% \\
Caldas & 36,6 & 15,6 & 2,36\% & 2,81\% \\
Sucre & 22,0 & 10,4 & 2,03\% & 2,79\% \\
San Andrés y Providencia & 82,4 & 33,3 & 2,86\% & 2,78\% \\
Bogotá D.C. & 67,4 & 28,6 & 2,05\% & 2,44\% \\
Boyacá & 47,5 & 22,6 & 1,49\% & 2,16\% \\
Tolima & 40,0 & 16,9 & 1,76\% & 1,75\% \\
Valle del Cauca & 54,5 & 23,5 & 1,39\% & 1,72\% \\
Santander & 62,0 & 26,1 & 1,40\% & 1,68\% \\
Chocó & 28,6 & 11,9 & 1,96\% & 0,89\% \\
Bolívar & 45,1 & 19,8 & 0,61\% & 0,82\% \\
Norte de Santander & 24,3 & 10,0 & 0,94\% & 0,79\% \\
Vichada & 92,2 & 38,1 & 1,33\% & 0,77\% \\
Nariño & 17,8 & 9,1 & -0,49\% & 0,46\% \\
Cauca & 26,6 & 13,7 & -0,21\% & 0,32\% \\
Córdoba & 22,0 & 10,7 & 0,16\% & 0,25\% \\
Antioquia & 47,7 & 20,3 & -0,09\% & 0,21\% \\
Amazonas & 50,2 & 21,8 & -0,09\% & 0,20\% \\
Magdalena & 23,9 & 10,3 & 0,00\% & 0,17\% \\
Guaviare & 39,1 & 17,0 & -1,00\% & -0,14\% \\
La Guajira & 24,8 & 11,1 & -0,71\% & -0,39\% \\
Arauca & 171,2 & 66,8 & -0,28\% & -0,84\% \\
Huila & 31,6 & 14,2 & -1,31\% & -0,86\% \\
Cesar & 32,6 & 13,8 & -1,61\% & -1,07\% \\
Guainía & 38,8 & 17,0 & -1,46\% & -1,23\% \\
Atlántico & 37,3 & 16,0 & -1,77\% & -1,33\% \\
Cundinamarca & 40,4 & 17,4 & -1,69\% & -1,53\% \\
Meta & 63,1 & 25,9 & -2,67\% & -2,07\% \\
Vaupés & 58,5 & 26,2 & -3,31\% & -2,75\% \\
Casanare & 165,2 & 67,8 & -3,31\% & -3,33\% \\
\bottomrule
\end{tabular}
\caption*{\footnotesize Nota: PIB por trabajador en millones de pesos constantes de 2015; PIB por hora en miles de pesos constantes de 2015. Bogotá se trata como departamento. Se excluye 2020 por no contar con GEIH anual comparable en la base del proyecto. Fuente: cálculos propios con DANE, Cuentas Nacionales Departamentales, y GEIH.}
\end{table}